\documentclass[
  aps,
  pra,
  reprint,
  nofootinbib
]{revtex4-2}

\usepackage[T1]{fontenc}
\usepackage{mathtools}
\usepackage{amssymb}
\usepackage{amsthm}
\usepackage{braket}
\usepackage{hyperref}
\usepackage{orcidlink}

\hypersetup{
  colorlinks=true,
  linkcolor=blue,
  citecolor=blue,
  urlcolor=blue,
  pdftitle={Answer Partitions, Oracle Response Spectra, and Exact Quantum Query Advantage},
  pdfauthor={Karl Svozil}
}

\newtheorem{theorem}{Theorem}
\newtheorem{proposition}{Proposition}
\newtheorem{corollary}{Corollary}

\begin{document}

\title{Answer Partitions, Oracle Response Spectra, and Exact Quantum Query Advantage}

\author{Karl Svozil \orcidlink{0000-0001-6554-2802}}
\email{karl.svozil@tuwien.ac.at}
\affiliation{Institute for Theoretical Physics, TU Wien,
Wiedner Hauptstra{\ss}e 8--10/136, 1040 Vienna, Austria}

\date{\today}

\begin{abstract}
An oracle problem has two parts: an answer partition, which says which hidden
instances must receive the same answer, and an access model, which says what a
query returns. We show that the access model can change the exact quantum query
complexity in ways no classical query can detect. One query solves a task
exactly if and only if some probe makes the output states of different answer
cells orthogonal; we state this criterion together with a certificate that
proves when it fails. For an oracle that applies a fixed unitary \(V\), its
response, to the target whenever the queried bit is 1, one query solves a
dictator (a function equal to a single input bit) exactly if and only if \(0\)
lies in the convex hull of the spectrum of \(V\), and a two-bit parity,
including Deutsch's problem, if and only if \(-1\) is an eigenvalue of \(V\);
no other total Boolean function is one-query exact. Hence a qutrit
transposition and a cyclic shift, which return the same values on a target
prepared in 0, need one and two exact queries; with bounded error, one query to the
shift still succeeds with the optimal probability \(7/8\). An inaccessible copy
of \(k\) address bits raises the cost of identifying an \(n\)-bit
Bernstein--Vazirani string from one query to \(\min\{k+1,n\}\), although every
classical value stays the same. For the
standard oracle we reorganize the known exact complexities of the 35 balanced
three-bit partitions and classify all 6435 balanced four-bit partitions
exactly; rational dual certificates prove the 6040 four-bit lower bounds that
were known only from numerical semidefinite programming.
\end{abstract}

\maketitle

\section{Introduction}

In Deutsch's problem a hidden function maps \(x\in\{0,1\}\) to a bit \(b_x\),
and one must decide whether it is constant (\(b_0=b_1\)) or balanced
(\(b_0\neq b_1\)). To keep the query unitary, the oracle does not overwrite a
register with \(b_x\) but adds \(b_x\) to a target bit \(y\),
\(\ket{x,y}\mapsto\ket{x,y\oplus b_x}\); started at the blank value \(y=0\),
the target then shows \(b_x\). With this XOR oracle one quantum query
suffices, against two classical ones~\cite{deutsch,deutsch:92,cleve:98}.

Now take a three-level target \(y\in\{0,1,2\}\) and require, for a fair
comparison, the same classical readout as for the XOR oracle: started at the
blank value \(y=0\), the target shows \(b_x\). This requirement fixes only
what happens to \(y=0\); what the oracle does with \(y=1\) and \(y=2\) is
left open. When
\(b_x=1\), the cyclic shift \(y\mapsto y+1\bmod 3\) moves
\(0\to1\to2\to0\), whereas the transposition exchanges \(0\) and \(1\) and
leaves \(2\) alone. Classically the two are equally useful: whatever basis
state the target starts in, one query reveals at most the bit \(b_x\), so
both need two classical queries. Yet with the shift an exact answer needs two
quantum queries, with the transposition only one (Sec.~\ref{sec:spectrum}).
A quantum algorithm can also prepare the target in an eigenstate of the
oracle's response, which the oracle only multiplies by the corresponding
eigenvalue. Deutsch's one-query solution needs the eigenvalue \(-1\). The bit
flip of the XOR oracle and the transposition both have it, with eigenstate
\(\ket0-\ket1\); the eigenvalues of the shift are the cube roots of unity, and
none of them is \(-1\). A classical query feeds in a single basis state and
never sees these eigenvalues.

In general, a query problem has two parts: an answer partition, which groups
the hidden instances into cells whose members share the correct answer, and an
access model, which says what each query returns. Query complexity, and with it
any quantum speedup, belongs to this pair, not to the task alone. That it
depends on the access model is known: simulating an in-place oracle for a
permutation of \(N\) points can require \(\Theta(\sqrt N)\) uses of the
standard oracle for the same function~\cite{KashefiKentVedralBanaszek2002};
the comparison of quantum and classical oracle power depends on how the oracle
is represented~\cite{WojcikChhajlany2006}; a richer classical oracle solves
Deutsch's problem with one query~\cite{Calude2006Dequantising}; and a faulty
oracle removes Grover's speedup~\cite{RegevSchiff2008}.

Earlier work, including the author's, posed quantum decision problems for
given quantum states: the possible states are partitioned into equivalence
classes, and a measurement must find the class of the state at hand;
identifying the state completely is the special case of one-element
classes~\cite{DonSvo01,svozil-2002-statepart-prl,svozil-2003-garda,svozil-2005-ko,2007-tkadlec-svozil-springer,svozil-2016-vector}.
Query problems are partitioned in the same way, but no state is given. The
hidden instance is a transformation that the algorithm can only call, an
oracle or black box. ``Oracle'' is the computer-science term, going back to
Turing's oracle machines~\cite{turing-1939}: something one puts questions to
and that answers. ``Black box'' stresses that one can only call it, not look
inside. We use ``oracle'' throughout. The algorithm must first produce states
by querying it. It succeeds exactly when the states produced for instances with
different answers are orthogonal, so that a measurement can tell them apart.
This paper asks which properties of the oracle, beyond the classical values it
returns, decide how many queries this takes.

The procedure has four steps, used throughout the paper and listed
in Sec.~\ref{sec:discussion}. First, the task fixes the answer partition:
instances with the same correct answer share a cell, and only the cell must be
learned. Second, the access model fixes what one oracle call does. Third, an
instance-independent probe, typically a superposition fanned out from a single
basis state, is sent through the oracle, which turns it into an output state
that depends on the hidden instance. Fourth, one call solves the task exactly
if and only if some probe makes the output states of different cells
orthogonal; overlaps within a cell do no harm. In the example above, this
orthogonality needs a relative factor \(-1\) between interfering branches,
which the bit flip and the transposition supply and the cyclic shift does not.

The main new result is a spectral criterion (Theorem~\ref{thm:spectrum} and
Proposition~\ref{prop:oneclass}). It concerns functions \(f(b)\) of a hidden
truth table \(b=(b_1,\ldots,b_N)\), in particular dictators and parities. A
dictator is a function whose value is just one of its input bits,
\(f(b)=b_i\) (or its negation): one bit dictates the answer and the others do
not matter, so classically one query, to address \(i\), suffices. For an
oracle that applies a fixed unitary \(V\), its response, to the target
whenever the queried bit is 1, one query solves a dictator exactly if
and only if \(0\) lies in the convex hull of the spectrum of \(V\), and a
two-bit parity \(b_i\oplus b_j\), including Deutsch's problem, if and only if
\(-1\) is an
eigenvalue of \(V\); no other total Boolean function is one-query exact. For
the XOR oracle this reduces to the characterization by Chen, Ye, and
Li~\cite{ChenYeLi2020}. It follows that two reversible oracles that agree on
the blank value \(y=0\) can have different exact costs, that a phase invisible
to basis queries can restore the advantage, and that a cyclic shift on an odd
number of levels keeps a bounded-error advantage although it loses the exact
one
(Sec.~\ref{sec:spectrum}). The second new result is an exact classification of
all balanced four-bit tasks for the standard oracle: rational dual
certificates prove the two-query lower bounds that were previously known only
numerically~\cite{MontanaroJozsaMitchison2015} (Sec.~\ref{sec:census}).
Section~\ref{sec:criterion} sets up the framework, a block form of standard
discrimination conditions and its dual. Section~\ref{sec:witness} treats an
oracle that leaves an inaccessible copy of the address; this changes
Bernstein--Vazirani costs although every classical value stays the same.
Table~\ref{tab:prior} separates prior results from new ones.

\begin{table*}[t]
\caption{Relation to prior work. ``Known'' lists results available before
this paper; ``This work'' lists what is added here. \(Q_E\) is the exact
quantum query complexity, \(\operatorname{Par}_N\) the parity of \(N\) bits,
\(V\) the response unitary, \(\operatorname{conv}(\operatorname{spec}V)\) the
convex hull of its spectrum, \(d\) the dimension of the cyclic response
\(X_d\ket y=\ket{y+1\bmod d}\), and \(k\) the number of the \(n\) address bits
recorded by the witness.}
\label{tab:prior}
\begin{ruledtabular}
\begin{tabular}{p{0.2\textwidth}p{0.36\textwidth}p{0.36\textwidth}}
Topic & Known & This work\\
\hline
Exact one-query criterion &
Perfect discrimination of unitaries~\cite{Acin2001UnitaryDiscrimination,DuanFengYing2009} &
Block form for answer partitions and its dual certificate
(Props.~\ref{prop:block} and \ref{prop:dual}); standard convex duality,
stated for use below\\
One-query exact class &
XOR oracle: dictators and two-bit parities~\cite{ChenYeLi2020};
partial functions~\cite{YeLi2020partial} &
Every response unitary \(V\): dictators iff
\(0\in\operatorname{conv}(\operatorname{spec}V)\), two-bit parities iff
\(-1\in\operatorname{spec}V\), no other total Boolean functions
(Thm.~\ref{thm:spectrum},
Prop.~\ref{prop:oneclass})\\
Dependence on the oracle model &
Oracle representation~\cite{WojcikChhajlany2006}; in-place versus standard
oracles~\cite{KashefiKentVedralBanaszek2002}; de-quantised Deutsch
oracle~\cite{Calude2006Dequantising}; faulty
oracles~\cite{RegevSchiff2008} &
A spectrum invisible to classical queries decides the exact cost: two
reversible oracles that agree on \(y=0\), shift with an added phase
(Cor.~\ref{cor:completions}); sign oracles cannot separate complementary truth
tables (Prop.~\ref{prop:sign})\\
Bounded error, odd cyclic response &
--- &
Optimal unambiguous one-query solution, worst-case success
\((3+\cos\frac{\pi}{d})/4\) (Prop.~\ref{prop:bounded})\\
Parity &
\(Q_E(\operatorname{Par}_N)=\lceil N/2\rceil\)~\cite{bbcmw-01,Farhi-98} &
Degree bound for every response; open qutrit question
(Sec.~\ref{sec:spectrum})\\
Three-bit functions &
Explicit optimal exact algorithms for all
functions~\cite{MontanaroJozsaMitchison2015} &
No new values; reorganization as balanced partitions
(Table~\ref{tab:n3})\\
Four-bit functions &
Numerical SDP values for all functions~\cite{MontanaroJozsaMitchison2015},
used in structural work~\cite{YeYangHuang2024} &
Exact rational dual certificates for all 6435 balanced partitions
(Table~\ref{tab:n4})\\
Decohered access &
Which-way information~\cite{Englert1996}; coherence in
Bernstein--Vazirani~\cite{NaseriEtAl2022} &
Exact costs with an inaccessible witness: surviving partition in one query
versus \(n-k\); identification in exactly \(\min\{k+1,n\}\) queries
(Prop.~\ref{prop:witness})\\
\end{tabular}
\end{ruledtabular}
\end{table*}

These results also correct a common picture of quantum parallelism: what
matters for the separations below is how the oracle responds (its phases,
spectrum, accessible registers, and environment), not how many branches a
superposition has.

\section{Answer partitions and exact one-query resolution}
\label{sec:criterion}

Let \(\Omega\) be a finite set of instances and let \(g:\Omega\rightarrow J\)
give the correct answer. The task defines the \emph{answer partition}
\begin{equation}
 \Pi_g=\{C_j:j\in g(\Omega)\},\qquad C_j=g^{-1}(j).
 \label{eq:task-partition}
\end{equation}
The partition says what must be found, not how. In the query model the hidden
instance labels an oracle; the algorithm is not given a state
\(\ket a\) that it could simply measure. The question is whether the given
oracle and a fixed protocol can bring the cell label into a register that can
be measured.

Let \(\rho_a^{(T)}\) be the state of the measurable registers after \(T\)
queries on instance \(a\). Exact resolution means that one instance-independent
measurement \(\{P_j\}\) satisfies
\begin{equation}
 P_jP_\ell=\delta_{j\ell}P_j,\quad
 \sum_jP_j=I,\quad
 \operatorname{Tr}(P_{g(a)}\rho_a^{(T)})=1.
 \label{eq:output-sectors}
\end{equation}
States from different cells then have orthogonal supports; states within a
cell need not be orthogonal. The cells \(C_j\) are sets of classical instances,
not subspaces; the subspaces \(P_j\mathcal H_{\rm out}\) of the output space
\(\mathcal H_{\rm out}\) exist only after an oracle and a protocol have
produced them. The partition thus prescribes which output overlaps must
vanish; it is not a ready-made observable.

For a partition \(\Pi\), a quantum oracle family \(\mathcal U\), and the
corresponding classical oracle \(\mathcal U_{\rm cl}\),
\(Q_\epsilon(\Pi;\mathcal U)\) and \(R_\epsilon(\Pi;\mathcal U_{\rm cl})\)
denote the smallest worst-case numbers of quantum and classical calls that
achieve worst-case error at most \(\epsilon\), with matching access on both
sides; \(Q_E\equiv Q_0\) is the exact quantum complexity, and \(D\) is the
worst-case number of queries of an exact deterministic classical algorithm
(the depth of an optimal decision tree). With worst-case counting, \(D=R_0\):
fixing the random choices of a zero-error randomized protocol that never uses
more than \(T\) queries leaves a deterministic protocol that is always correct
with at most \(T\) queries. (Counted in expectation, as is common, zero-error
cost can be smaller than \(D\).) A deterministic classical protocol groups the
instances by the transcript of answers they produce; it succeeds exactly only
if instances with the same transcript always lie in the same cell. A quantum
protocol keeps no such intermediate record; it changes the overlaps among the
instance states and measures only after the cells have separated. A speedup is
a gap between the quantum and the classical cost for the same task and the
same access.

All Hilbert spaces are finite dimensional. For a family of unitary oracles
\(\{U_a:a\in\Omega\}\), a general one-query protocol prepares a pure state
\(\ket\Psi\) of the query system \(Q\) and a reference system \(R\), applies
\(U_a\otimes I_R\), and then applies a fixed channel and measurement. Write
\begin{equation}
 \ket{\Psi_a}=(U_a\otimes I_R)\ket\Psi,
 \qquad
 \sigma=\operatorname{Tr}_R\ket\Psi\!\bra\Psi .
\end{equation}
The Gram matrix of the output states is
\begin{equation}
 G_{aa'}=\braket{\Psi_a|\Psi_{a'}}
 =\operatorname{Tr}\!\left(\sigma U_a^\dagger U_{a'}\right).
 \label{eq:gram}
\end{equation}

\begin{proposition}[Block discrimination]\label{prop:block}
One query to \(\{U_a\}\), with any reference system allowed, resolves
\(\Pi_g\) exactly if and only if some density operator \(\sigma\) satisfies
\begin{equation}
 \operatorname{Tr}\!\left(\sigma U_a^\dagger U_{a'}\right)=0
 \quad\text{whenever}\quad g(a)\ne g(a').
 \label{eq:criterion}
\end{equation}
\end{proposition}

\begin{proof}
Exact discrimination requires the output states of different cells to span
orthogonal subspaces, that is, Eq.~\eqref{eq:criterion} for every pair from
different cells; no fixed channel can make nonorthogonal states perfectly
distinguishable. Conversely, orthogonal subspaces are distinguished by the
projectors onto them. By purification, optimizing over \(\sigma\) covers mixed
probes and arbitrary ancillas.
\end{proof}

Equation~\eqref{eq:criterion} is the block form of the standard conditions for
perfect discrimination of unitaries
\cite{Acin2001UnitaryDiscrimination,DuanFengYing2009}; it is not a new theorem
about channel discrimination. Here it links an answer partition to a given
oracle family. The Gram matrix need only be block diagonal with respect to
\(\Pi_g\); states within one cell may overlap or coincide. Identifying the
instance completely is the special case of one-element cells, where \(G\) must
be diagonal. Since zero queries cannot resolve a nontrivial partition,
Eq.~\eqref{eq:criterion} characterizes \(Q_E(\Pi_g;\mathcal U)=1\). With
bounded error, orthogonality is replaced by a measurement that reaches the
required success probability. Lower bounds for several queries come from
adversary methods~\cite{Ambainis2002Adversary}, whose general form
characterizes state conversion and function
evaluation~\cite{LeeEtAl2011StateConversion,Reichardt2014SpanPrograms}.

When Eq.~\eqref{eq:criterion} has no solution, a single matrix proves it.

\begin{proposition}[Dual certificate]\label{prop:dual}
Let \(X=\{(a,a'):g(a)\neq g(a')\}\). Exactly one of the following holds:
(i) a density operator \(\sigma\) satisfies Eq.~\eqref{eq:criterion};
(ii) there is a complex matrix \(\Gamma\), supported on \(X\) and satisfying
\(\Gamma_{a'a}=\Gamma_{aa'}^*\), such that
\begin{equation}
 Z(\Gamma)=\sum_{(a,a')\in X}\Gamma_{aa'}\,U_a^\dagger U_{a'}\succ0 .
 \label{eq:dual-witness}
\end{equation}
\end{proposition}

\begin{proof}
\(Z(\Gamma)\) is Hermitian because \((U_a^\dagger U_{a'})^\dagger=
U_{a'}^\dagger U_a\). If (i) and (ii) both held, then
\(\operatorname{Tr}[\sigma Z(\Gamma)]=0\), contradicting \(Z(\Gamma)\succ0\).
If (ii) fails, the real subspace \(\{Z(\Gamma)\}\) of Hermitian operators misses
the open convex cone of positive definite operators, so a separating
hyperplane gives a nonzero \(\sigma\) with \(\operatorname{Tr}(\sigma P)\ge0\)
for all \(P\succ0\), that is \(\sigma\succeq0\), and
\(\operatorname{Tr}[\sigma Z(\Gamma)]=0\) for all admissible \(\Gamma\). Taking
\(\Gamma\) supported on a single pair \((a,a')\) and its transpose, with real
and with imaginary entries, gives
\(\operatorname{Tr}(\sigma U_a^\dagger U_{a'})=0\); normalizing \(\sigma\) gives
(i).
\end{proof}

Like an adversary matrix, \(\Gamma\) weights only pairs from different cells.
After \(t\) queries the same weights define the progress measure
\(\sum_{(a,a')\in X}\Gamma_{aa'}\braket{\psi_a^{(t)}|\psi_{a'}^{(t)}}\) of the
adversary method~\cite{Ambainis2002Adversary}, and the query cost measures how
fast such weighted cross-cell overlaps can be driven to zero; this is what it
means for a quantum algorithm to ``separate'' answer cells. When all
\(U_a^\dagger U_{a'}\) commute, as for the oracles of Secs.~\ref{sec:spectrum}
and~\ref{sec:census}, which are
diagonal in the eigenbasis of their response, only the diagonal of \(\sigma\)
in a common eigenbasis enters Eq.~\eqref{eq:criterion}, and the criterion
becomes a linear feasibility problem for a probability vector.

The criterion also separates a mere change of representation from a change of
access. If
\begin{equation}
 V_a=A U_a B
 \label{eq:fixed-equivalence}
\end{equation}
for fixed unitaries \(A,B\), then
\(V_a^\dagger V_{a'}=B^\dagger U_a^\dagger U_{a'}B\), so both families reach
the same Gram matrices. More generally, families that simulate each other
exactly with \(O(1)\) calls have the same quantum query complexity up to a
constant factor; an approximate simulation adds its error to \(\epsilon\). For
a quantum--classical gap to carry over, the classical oracles must simulate
each other in the same way. Families not related by
Eq.~\eqref{eq:fixed-equivalence} can differ in their exact one-query cost, as
the next section shows.

\section{The response spectrum decides exact one-query power}
\label{sec:spectrum}

\subsection{Deutsch's problem and the sign oracle}

Let \(b=(b_0,b_1)\) be the truth table of a Boolean function and use, with
address \(x\) and target bit \(y\),
\begin{equation}
 O_b\ket{x,y}=\ket{x,y\oplus b_x}.
 \label{eq:xor-oracle}
\end{equation}
With the address in \(\ket+\) and the target qubit in \(\ket-\), phase kickback
leaves the address register in the states \(\ket{\psi_b}\), labelled by the
truth table \(b=b_0b_1\) (the target qubit factors out),
\begin{equation}
 \ket{\psi_{00}} = -\ket{\psi_{11}} = \ket{+},\quad
 \ket{\psi_{01}} = -\ket{\psi_{10}} = \ket{-}.
 \label{eq:deutsch-rays}
\end{equation}
Constant and balanced functions give orthogonal states, while the two
functions within each cell give the same state up to sign. The Gram matrix is
block diagonal but not diagonal, and the result is the familiar separation of
one quantum versus two classical
queries~\cite{deutsch,deutsch:92,cleve:98}.

Hadamard gates on the target qubit turn Eq.~\eqref{eq:xor-oracle} into the
controlled phase \(\ket{x,z}\mapsto(-1)^{z b_x}\ket{x,z}\), so the two forms
are equivalent in the sense of Eq.~\eqref{eq:fixed-equivalence}. The sign
oracle \(P_b\ket x=(-1)^{b_x}\ket x\), without the control qubit \(z\), is
weaker. A basis query to it reveals nothing about \(b_x\), and for the
complemented truth table \(\bar b\) one has \(P_{\bar b}=-P_b\), so the two
differ only by a global phase.

\begin{proposition}[Complement obstruction]\label{prop:sign}
With access to \(\{P_b\}\) alone (any ancillas, but no control on \(P_b\)), no
number of queries resolves a partition that puts some \(b\) and its complement
\(\bar b\) in different cells, even with error \(\epsilon<1/2\).
\end{proposition}

\begin{proof}
Each use of \(P_{\bar b}\otimes I\) equals \(-(P_b\otimes I)\), and every branch
of an adaptive protocol picks up the same sign. The final density operators for
\(b\) and \(\bar b\) are therefore identical, and so are the output
distributions.
\end{proof}

Deutsch's constant--balanced partition is closed under complement and is still
solved by one query to \(P_b\). Three-bit majority, the parity of an odd number
of bits, and every other self-dual function cannot be resolved at all. Of the
balanced partitions in Sec.~\ref{sec:census}, only 3 of 35 for three bits (the
two-bit parities) and 35 of 6435 for four bits are closed under complement;
the sign oracle resolves none of the others, while the XOR oracle resolves all
of them with at most two (\(N=3\)) or three (\(N=4\)) queries. A coherent
control on \(P_b\) would restore the missing relative phase, but control is
extra access to the oracle, not something one can add to an unknown
unitary~\cite{AraujoFeixCostaBrukner2014}.

\subsection{Response unitaries}

Let \(V\) be any unitary and consider the family, indexed by the same four
truth tables, that applies \(V\) to the target whenever \(b_x=1\):
\begin{equation}
 U_b^{(V)}
 =\sum_{x=0}^1\ket x\!\bra x\otimes V^{b_x}.
 \label{eq:response-oracle}
\end{equation}
\emph{Access conventions.} \(V\) is known. Between queries the algorithm may
apply any fixed operation, including preparing eigenvectors of \(V\), at no
query cost; it may not control, invert, or otherwise modify the oracle unless
that is explicitly allowed. The phase of \(V\) matters: \(V^{b_x}\) acts only on
the branch selected by the address, so replacing \(V\) by \(e^{i\theta}V\)
multiplies Eq.~\eqref{eq:response-oracle} by \(e^{i\theta b_x}\), a phase that
differs between address branches and is therefore observable; it can change
the spectral condition below (Corollary~\ref{cor:completions}). This differs
from access to an unknown unitary or channel, which is defined only up to a
global phase.

\begin{theorem}[Response spectrum]\label{thm:spectrum}
One query to \(\{U_b^{(V)}\}\), with any probe and reference, solves Deutsch's
problem exactly if and only if \(-1\) is an eigenvalue of \(V\).
\end{theorem}

\begin{proof}
If \(V\ket w=-\ket w\), the probe \(\ket+\ket w\) gives ordinary phase kickback
and Eq.~\eqref{eq:deutsch-rays}. Conversely, give weight \(\Gamma=1\) to the
four cross-cell pairs (and their transposes) in
Proposition~\ref{prop:dual}. The pairs \((00,01)\), \((00,10)\), \((11,01)\),
and \((11,10)\) contribute
\(I\oplus V\), \(V\oplus I\), \(V^\dagger\oplus I\), and \(I\oplus V^\dagger\)
on the address blocks \(x=0,1\). Together with their adjoints, each address
block receives \(4I+2V+2V^\dagger\), and since \(V\) is unitary,
\begin{equation}
 Z=2\,I_2\otimes(2I+V+V^\dagger)
 =2\,I_2\otimes(I+V)(I+V)^\dagger .
 \label{eq:deutsch-witness}
\end{equation}
\(Z\) is positive definite unless \(I+V\) is singular, that is, unless \(-1\) is
an eigenvalue of \(V\).
\end{proof}

In words: at each of the two addresses the probe sees an average of
eigenvalues of \(V\), a point of \(\operatorname{conv}(\operatorname{spec}V)\),
and orthogonality across cells forces both points to \(-1\). Since \(-1\) lies
on the unit circle, an average of eigenvalues equals \(-1\) only if \(-1\) is
itself an eigenvalue. Qudit versions of Deutsch--Jozsa enlarge the address and
response systems~\cite{Cereceda2004QuditDeutsch,deutsch:92,cleve:98};
Theorem~\ref{thm:spectrum} instead keeps the four Deutsch instances fixed and
varies the response unitary. It also differs from ``useless-query'' bounds,
which derive quantum lower bounds, under a prior distribution, from classical
queries that reveal nothing about a property~\cite{MeyerPommersheim2011}.

\begin{corollary}[Same classical values, different cost]\label{cor:completions}
On a qutrit, consider the transposition \(\ket0\leftrightarrow\ket1\) that
fixes \(\ket2\), the cyclic shift \(X_3\ket y=\ket{y+1\bmod 3}\), and
\(X_3'=e^{i\pi/3}X_3\). All three map the blank target \(\ket0\) to \(\ket1\)
up to a phase (\(X_3'\ket0=e^{i\pi/3}\ket1\)), so measuring a blank target
returns \(b_x\) in each case; \(X_3\) and \(X_3'\)
moreover agree on every computational-basis input up to a phase. Their
spectra are \(\{1,1,-1\}\), \(\{1,e^{2\pi i/3},e^{4\pi i/3}\}\), and
\(\{e^{i\pi/3},-1,e^{-i\pi/3}\}\), so the exact quantum cost of Deutsch's
problem is \(Q_E=1\), \(2\), and \(1\), respectively, while \(D=2\) in all three
cases. More generally, for \(X_d\ket y=\ket{y+1\bmod d}\) with \(d\ge2\),
\begin{equation}
 Q_E^{(d)}=
 \begin{cases}1,&d\ \text{even},\\2,&d\ \text{odd},\end{cases}
 \qquad D^{(d)}=2.
 \label{eq:cyclic-response-cost}
\end{equation}
\end{corollary}

Specifying an oracle by what it does to a blank target,
\(\ket{x}\ket{0}\mapsto\ket{x}\ket{b_x}\), fixes only one column of each
\(V^{b_x}\). How the rest is filled in, by a transposition or a cyclic shift,
decides whether Deutsch's exact advantage exists, and so does a phase that no
basis query can see: the factor \(e^{i\pi b_x/3}\) in \(X_3'^{\,b_x}\) differs
between address branches; it is a phase oracle composed with the shift.
Neither the classical values nor even the full action on basis states up to
phases determines \(Q_E\). The effect is a constant factor, not an asymptotic
separation.

\subsection{Any number of bits}

The same argument works for any number of bits. Let
\(f:\{0,1\}^N\to\{0,1\}\) be nonconstant, with oracle
\(U_b^{(V)}=\sum_{i=1}^N\ket i\!\bra i\otimes V^{b_i}\). Below, \(e_i\) is the
truth table with a single 1 at address \(i\), and \(0^N\) the all-zero truth
table.

\begin{proposition}[One-query class]\label{prop:oneclass}
One query to \(\{U_b^{(V)}\}\) resolves a nonconstant \(f\) exactly if and only
if, up to negation, either \(f=b_i\) and
\(0\in\operatorname{conv}(\operatorname{spec}V)\), or \(f=b_i\oplus b_j\) and
\(-1\in\operatorname{spec}V\). The dictator condition
\(0\in\operatorname{conv}(\operatorname{spec}V)\) holds in particular if
\(V\ket y\perp\ket y\) for some basis state \(\ket y\).
\end{proposition}

\begin{proof}
Let \(\rho_i\succeq0\) be the diagonal block of \(\sigma\) at address \(i\), an
operator on the target, and \(p_i=\operatorname{Tr}\rho_i\), so
\(\sum_ip_i=1\). Since
\(U_b^\dagger U_{b'}=\sum_i\ket i\!\bra i\otimes V^{b'_i-b_i}\), only the
\(\rho_i\) enter Eq.~\eqref{eq:criterion}. If \(f\) depends on \(b_i\), some
\(b\) has \(f(b)\neq f(b\oplus e_i)\), and Eq.~\eqref{eq:criterion} reads
\(1-p_i+\operatorname{Tr}(\rho_iV^{\pm1})=0\). Since
\(|\operatorname{Tr}(\rho_iV^{\pm1})|\le p_i\), this forces \(p_i\ge1/2\), so
\(f\) depends on at most two bits. If it depends on two, \(i\) and \(j\), then
\(p_i=p_j=1/2\) and \(\operatorname{Tr}(\rho_iV^{\pm1})=-p_i\). An average of
unit-modulus eigenvalues equals \(-1\) only if all its weight lies on the
eigenvalue \(-1\), so \(\rho_i\) and \(\rho_j\) are supported on the \(-1\)
eigenspace of \(V\). The overlap of the
outputs for \(b\) and \(b\oplus e_i\oplus e_j\) is then \(-1\neq0\), so these
instances share a cell, \(f(b)=f(b\oplus e_i\oplus e_j)\), and \(f\) is
\(b_i\oplus b_j\) or its negation. A \(-1\) eigenvector gives the Deutsch
construction on addresses \(i,j\). For a dictator \(b_i\), the cross-cell
pairs \((0^N,e_i)\) and \((0^N,e_i\oplus e_j)\), \(j\neq i\), give
\(\operatorname{Tr}(\rho_iV)+\sum_{k\neq i}p_k=0\) and, after subtracting,
\(\operatorname{Tr}(\rho_jV)=p_j\) for every \(j\neq i\). The density operator
\(\tau=\sum_k\rho_k\) therefore has \(\operatorname{Tr}(\tau V)=0\), that is,
\(0\in\operatorname{conv}(\operatorname{spec}V)\). Conversely, the probe
\(\ket i\ket w\) with \(\braket{w|V|w}=0\) makes every cross-cell overlap
vanish. With an additional null address, on which the oracle acts as the
identity (as in the phase form used in Sec.~\ref{sec:census}), the same
argument holds with \(\operatorname{spec}V\) replaced by
\(\operatorname{spec}V\cup\{1\}\) in the dictator condition.
\end{proof}

For a dictator \(b_i\) it suffices to query address \(i\) alone; one query
then succeeds exactly when some target state \(\ket w\) is turned into an
orthogonal one, \(\braket{w|V|w}=0\), that is, when an average of eigenvalues
of \(V\) vanishes. For \(V=-I\), the family is the sign oracle of
Proposition~\ref{prop:sign}:
one query resolves two-bit parities, but no number of queries resolves a
dictator.

For the XOR oracle, \(V=X\), Proposition~\ref{prop:oneclass} reproduces the
characterization of exact one-query algorithms by Chen, Ye, and
Li~\cite{ChenYeLi2020}. Proposition~\ref{prop:oneclass} concerns total
functions. With a promise, more tasks can be one-query exact, the
Deutsch--Jozsa problem~\cite{deutsch:92} being the standard example; partial
functions are treated in Ref.~\cite{YeLi2020partial}. For an odd cyclic
response \(X_d\) only the dictators remain, and a single classical query
already resolves them.

\subsection{Bounded error for odd cyclic responses}

Corollary~\ref{cor:completions} concerns exact solutions. With bounded error,
the odd cyclic response keeps a quantum advantage. As a referee pointed out,
a single eigenvector of \(X_3\) with eigenvalue \(e^{2\pi i/3}\), the address in \(\ket+\), and a
measurement of the address in \(\{\ket+,\ket-\}\) already give success 1 on
constant and \(3/4\) on balanced instances, whereas one classical query cannot
exceed worst-case success \(1/2\). Spreading the probe over the two
eigenvalues adjacent to \(-1\) raises the worst case further and makes every
conclusive answer correct.

\begin{proposition}[Optimal one-query solution]\label{prop:bounded}
For odd \(d\ge3\), one query to \(\{U_b^{(X_d)}\}\) solves Deutsch's problem
with zero error and a conclusive answer with probability
\((1+\cos\frac{\pi}{d})/2\) on every instance. Guessing on the inconclusive
outcome gives worst-case success \((3+\cos\frac{\pi}{d})/4\), that is \(3/4\)
conclusive and \(7/8\) overall for \(d=3\), and no one-query protocol does
better. One classical query never gives a conclusive zero-error answer and has
worst-case success \(1/2\).
\end{proposition}

\begin{proof}
Let \(\ket{v_\pm}\) be eigenvectors of \(X_d\) with the eigenvalues
\(\lambda_\pm=-e^{\pm i\pi/d}\) adjacent to \(-1\), and use the probe
\(\frac12\sum_{x=0}^1\sum_\pm\ket x\ket{v_\pm}\). Writing
\(\gamma=\cos\frac{\pi}{d}\) and using
\((\lambda_+^{k}+\lambda_-^{k})/2=-\gamma\) for \(k=\pm1\), the output Gram
matrix is \(G_{bb'}=\frac12\sum_x c_x\) with \(c_x=1\) if \(b_x=b'_x\) and
\(c_x=-\gamma\) otherwise. Hence
\(\ket{\psi_{00}}+\ket{\psi_{11}}=\ket{\psi_{01}}+\ket{\psi_{10}}\), with squared
norm \(2(1-\gamma)\), while \(\ket{\psi_{00}}-\ket{\psi_{11}}\) and
\(\ket{\psi_{01}}-\ket{\psi_{10}}\) are orthogonal to each other and to that
common vector, each with squared norm \(2(1+\gamma)\). Projecting onto these
three directions returns ``constant,'' ``balanced,'' or ``inconclusive''; each
instance gives the correct conclusive outcome with probability
\((1+\gamma)/2\) and never the wrong one.

For optimality, write any pure probe, including a reference, as
\(\ket0\ket A+\ket1\ket B\) with \(\lVert A\rVert^2+\lVert B\rVert^2=1\), and
put \(W=V\otimes I_R\) with \(V=X_d\). Let \(\rho_C\) and \(\rho_B\) be the
outputs averaged uniformly over the constant and over the balanced instances.
Their diagonal address blocks coincide, and their off-diagonal blocks differ
by \(\frac12\ket{(I-W)A}\!\bra{(I-W)B}\), so
\(\lVert\rho_C-\rho_B\rVert_1=\lVert(I-W)A\rVert\,\lVert(I-W)B\rVert
\le\frac12\lVert I-V\rVert^2\), where \(\lVert\cdot\rVert_1\) is the trace norm
and \(\lVert\cdot\rVert\) the operator norm (for vectors, the Euclidean norm).
Worst-case success cannot exceed success under
this prior, which by the Helstrom bound is at most
\(\frac12+\frac18\lVert I-V\rVert^2\). For odd \(d\),
\(\lVert I-X_d\rVert^2=2+2\cos\frac{\pi}{d}\), which gives
\((3+\cos\frac{\pi}{d})/4\). Classically, one query reveals one bit \(b_x\),
and both cells remain consistent with every transcript.
\end{proof}

The bound also makes the construction optimal among unambiguous protocols. If
every instance gave a conclusive answer with probability at least \(p\),
guessing on the inconclusive outcome would give worst-case success at least
\((1+p)/2\); hence \(p\le(1+\cos\frac{\pi}{d})/2\), which the construction
attains.

The bound \(\frac12+\frac18\lVert I-V\rVert^2\) holds for every response \(V\)
and reaches 1 only if \(\lVert I-V\rVert=2\), that is, only if \(-1\) is an
eigenvalue; this is a second proof of the necessity part of
Theorem~\ref{thm:spectrum}. The construction parallels
Deutsch's original algorithm, which also answered conclusively only part of
the time~\cite{deutsch}; here the inconclusive probability
\((1-\cos\frac{\pi}{d})/2\) is set by how close the spectrum comes to \(-1\),
and it vanishes as \(d\to\infty\).

\subsection{Parity}

Proposition~\ref{prop:oneclass} shows that one query resolves the parity
\(\operatorname{Par}_N(b)=b_1\oplus\cdots\oplus b_N\) only for \(N\le2\). Proposition~\ref{prop:dual} gives an explicit
witness for the XOR oracle: with \(\Gamma=1\) on the pairs \((0^N,e_i)\),
\begin{equation}
 Z=\sum_{i=1}^N\bigl(O_{e_i}+O_{e_i}^\dagger\bigr)=2N\,I-4\Pi_-,
 \label{eq:parity-witness}
\end{equation}
where \(\Pi_-=\sum_{i}\ket i\!\bra i\otimes\ket-\!\bra-\); \(Z\) is positive
definite exactly when \(N>2\). For more queries, note that
\(U_b^{(V)}=\sum_i\ket i\!\bra i\otimes[I+b_i(V-I)]\) is affine in \(b\) for
every response \(V\). The amplitudes after \(T\) queries are therefore
polynomials of degree at most \(T\) in the input bits, the acceptance
probability has degree at most \(2T\), and parity has exact degree \(N\). Thus
\(T\ge N/2\) for every response, and for the XOR oracle applying the Deutsch
construction to disjoint pairs gives the matching algorithm, so
\begin{equation}
 Q_E(\operatorname{Par}_N)=\lceil N/2\rceil,
 \qquad D(\operatorname{Par}_N)=N
 \label{eq:parity-complexity}
\end{equation}
for the XOR oracle~\cite{bbcmw-01,Farhi-98}. For the qutrit shift \(X_3\),
Theorem~\ref{thm:spectrum} forbids the pair construction, and
\(\lceil N/2\rceil\le Q_E(\operatorname{Par}_N;X_3)\le N\). Whether the upper
bound is attained is open; if it is, two oracles with identical classical
values would differ by a factor of two for all \(N\).

If instead an oracle
\(\widetilde O_b\ket y=\ket{y\oplus\operatorname{Par}(b)}\) that returns the
parity counts as one call, parity takes one query, classically and quantum
mechanically. This does not contradict Eq.~\eqref{eq:parity-complexity}; such
an oracle already returns the answer.

\section{Exact census of balanced tasks for the standard oracle}
\label{sec:census}

So far the access varied and the task stayed fixed. Now the standard oracle
stays fixed and the task varies. The hidden instance is \(b\in\{0,1\}^N\),
accessed through Eq.~\eqref{eq:xor-oracle} with \(N\) addresses, and we
consider every two-valued answer with \(2^{N-1}\) instances in each cell.
Swapping the two answers gives the same partition, so the number of such
partitions is
\begin{equation}
 M_N=\frac12\binom{2^N}{2^{N-1}}.
 \label{eq:balanced-count}
\end{equation}
Permuting coordinates, complementing coordinates, and swapping the answers
preserve both \(D\) and \(Q_E\). These symmetries reduce the \(M_3=35\)
partitions to 6 orbits and the \(M_4=6435\) partitions to 58. The counts below
include the orbit sizes.

\emph{Methods.} \(D\) was computed exactly by dynamic programming over
subcubes. By Proposition~\ref{prop:oneclass} (Ref.~\cite{ChenYeLi2020}),
\(Q_E=1\) exactly for dictators and two-bit parities. Upper bounds on \(Q_E\)
come from adaptive trees whose nodes read either one bit or the parity of two
bits, at one query each; such trees prove upper bounds but do not capture all
exact quantum algorithms~\cite{MontanaroJozsaMitchison2015}. Lower bounds
\(Q_E\ge3\) come from the two-query semidefinite program of Barnum, Saks, and
Szegedy~\cite{BarnumSaksSzegedy2003}, in its bounded-error form, which is
numerically better behaved than an exact feasibility test. In the phase form of the
standard oracle, \(\ket i\mapsto(-1)^{b_i}\ket i\) with a null address \(i=0\) and
\(b_0=0\), put \(D_i=\operatorname{diag}_b\bigl((-1)^{b_i}\bigr)\) (so
\(D_0=I\)). A two-query algorithm with worst-case success \(p\) yields
positive semidefinite \(2^N\times2^N\) matrices \(A_i\), \(B_i\), \(\Gamma_0\),
\(\Gamma_1\) with \(\sum_iA_i=\mathbf 1\mathbf 1^{\mathsf T}\), where
\(\mathbf 1\) is the all-ones vector, \(\sum_iD_iA_iD_i=\sum_iB_i\),
\(\sum_iD_iB_iD_i=\Gamma_0+\Gamma_1\), and \((\Gamma_{f(b)})_{bb}\ge p\). If
real symmetric \(Y_1,Y_2,Y_3\) and a probability vector \(\lambda\) satisfy
\begin{equation}
\begin{aligned}
 &Y_1\succeq D_iY_2D_i,\quad Y_2\succeq D_iY_3D_i\quad(i=0,\dots,N),\\
 &Y_3\succeq\Lambda_j\quad(j=0,1),
\end{aligned}
 \label{eq:sdp-dual}
\end{equation}
with \(\Lambda_j=\operatorname{diag}_b\bigl(\lambda_b[f(b)=j]\bigr)\), where
\([\cdot]\) is 1 if the condition holds and 0 otherwise, then, with
\(\langle Y,A\rangle=\operatorname{Tr}(YA)\),
\begin{equation}
\begin{aligned}
 \mathbf 1^{\mathsf T}Y_1\mathbf 1&=\textstyle\sum_i\langle Y_1,A_i\rangle
 \ge\langle Y_2,\textstyle\sum_iB_i\rangle
 \ge\textstyle\sum_j\langle Y_3,\Gamma_j\rangle\\
 &\ge\textstyle\sum_b\lambda_b(\Gamma_{f(b)})_{bb}\ge p ,
\end{aligned}
 \label{eq:weak-duality}
\end{equation}
so \(\mathbf 1^{\mathsf T}Y_1\mathbf 1<1\) proves \(Q_E\ge3\). We solved
Eq.~\eqref{eq:sdp-dual} numerically with a barrier method, added
\(3\varepsilon\), \(2\varepsilon\), and \(\varepsilon\) times the identity to
\(Y_1\), \(Y_2\), and \(Y_3\) with \(\varepsilon=2^{-20}\), which raises every
constraint matrix in Eq.~\eqref{eq:sdp-dual} by \(\varepsilon I\), rounded
every entry to a dyadic rational with denominator \(2^{40}\), and verified all
twelve matrix inequalities and \(\lambda\ge0\), \(\sum_b\lambda_b=1\) in exact
integer arithmetic (positive pivots in Gaussian
elimination). Each certificate is a list of integers, given in the
Supplemental Material, and can be checked without the solver.

\emph{Use of AI tools.} The census code and the Lean formalization were
written with an OpenAI Codex assistant, and the dual-certificate solver and
the exact verifier with Anthropic's Claude (versions are given in the
Acknowledgments). No
reported result rests on these tools alone: the census counts were reproduced
by two independently written programs and agree with
Ref.~\cite{MontanaroJozsaMitchison2015}, the Lean audits reproduce the finite
census, and every certificate is checked by a verifier that uses only exact
integer arithmetic and does not depend on the solver.

\begin{table}[t]
\caption{Balanced three-bit tasks up to coordinate permutations, input
complements, and answer exchange. The orbit sizes sum to 35. Here
\(\operatorname{Par}_3\) and \(\operatorname{Maj}_3\) are the parity and the
majority of the three bits, and \(\mathrm{Sel}(c,u,v)=u\) for \(c=0\) and \(v\)
for \(c=1\). The values of
\(Q_E\) are not new: Ref.~\cite{MontanaroJozsaMitchison2015} gives explicit
optimal exact algorithms for every function on up to three bits.}
\label{tab:n3}
\begin{ruledtabular}
\begin{tabular}{lrrr}
Representative & Orbit & \(D\) & \(Q_E\)\\
\hline
\(b_1\) & 3 & 1 & 1\\
\(b_1\mathbin\oplus b_2\) & 3 & 2 & 1\\
\(\operatorname{Par}_3\) & 1 & 3 & 2\\
\(\operatorname{Maj}_3\) & 4 & 3 & 2\\
\(b_1\mathbin\oplus b_2b_3\) & 12 & 3 & 2\\
\(\mathrm{Sel}(b_1,b_2,b_3)\) & 12 & 2 & 2\\
\end{tabular}
\end{ruledtabular}
\end{table}

\begin{table*}[t]
\caption{Balanced four-bit tasks grouped by \((D,Q_E)\). Orbits are taken up
to coordinate permutations, input complements, and answer exchange; the
columns sum to 58 orbits and 6435 partitions. The last column names the source
of the lower bound on \(Q_E\): Proposition~\ref{prop:oneclass} excludes one
query, and the rational dual certificates of Eq.~\eqref{eq:sdp-dual} exclude
two.}
\label{tab:n4}
\begin{ruledtabular}
\begin{tabular}{crrll}
\((D,Q_E)\) & Orbits & Partitions & Representative & Lower bound\\
\hline
\((1,1)\) & 1 & 4 & \(b_1\) & ---\\
\((2,1)\) & 1 & 6 & \(b_1\oplus b_2\) & ---\\
\((2,2)\) & 1 & 48 & \(\mathrm{Sel}(b_1,b_2,b_3)\) & Prop.~\ref{prop:oneclass}\\
\((3,2)\) & 8 & 248 & \(\operatorname{Maj}(b_1,b_2,b_3)\) & Prop.~\ref{prop:oneclass}\\
\((4,2)\) & 4 & 89 & \(\operatorname{Par}_4\) & Prop.~\ref{prop:oneclass}\\
\((3,3)\) & 10 & 1536 & \(\mathrm{Sel}(b_1,b_2b_3,b_2\lor b_4)\) & certificate\\
\((4,3)\) & 33 & 4504 & \(b_1\oplus b_2\oplus b_3b_4\) & certificate\\
\end{tabular}
\end{ruledtabular}
\end{table*}

\emph{Three bits.} For \(N=3\) the lower bounds from the one-query class match
the explicit trees for every task (Table~\ref{tab:n3}). These values are known:
Montanaro, Jozsa, and Mitchison give optimal exact algorithms for every
three-bit function~\cite{MontanaroJozsaMitchison2015}; we only arrange them as
balanced partitions. Of the 35 balanced tasks, 20 have an exact advantage:
three by a factor of two and seventeen by a factor of \(3/2\). The trees show
how such algorithms work. For three-bit majority, first obtain
\(b_1\oplus b_2\); if it is 0, read \(b_1\), otherwise read \(b_3\); this is
the algorithm of Ref.~\cite{MontanaroJozsaMitchison2015}, Sec.~6.2. Likewise,
for \(F=b_1\oplus b_2b_3\), first read \(b_2\), then \(b_1\) if \(b_2=0\), or
the Deutsch parity \(b_1\oplus b_3\) if \(b_2=1\).

\emph{Four bits.} For \(N=4\) the one-query class and explicit two-query trees
settle 15 orbits covering 395 partitions. For each of the remaining 43 orbit
representatives, covering 6040 partitions, explicit trees give \(Q_E\le3\) and
the exact rational certificates bound the optimal two-query worst-case success
by values between 0.9001 and 0.9862, hence \(Q_E=3\). The classification of Table~\ref{tab:n4} is
therefore exact: 4847 balanced four-bit partitions have an exact quantum
advantage (6 by a factor of two at \((2,1)\), 89 by a factor of two at
\((4,2)\), 248 by a factor of \(3/2\), and 4504 by a factor of \(4/3\)) and
1588 have none. The polynomial method cannot give these lower bounds: the 43
orbits have real degree 3 (2144 partitions) or 4 (3896 partitions), so it
yields only \(Q_E\ge2\).

Ref.~\cite{MontanaroJozsaMitchison2015} reported numerical SDP values for all
Boolean functions of up to four bits and noted that such values do not by
themselves prove exactness; later work on exact two-query algorithms builds on
that classification~\cite{YeYangHuang2024}. We checked our orbit
representatives against the published NPN (negation--permutation--negation)
identifiers of Ref.~\cite{MontanaroJozsaMitchison2015}. After removing dummy
variables, its one- and two-query entries agree with all 15 representatives
settled by trees, and for the other 43 its optimal two-query success
probabilities, 0.900 to 0.986, agree with the optimal values 0.9000--0.9861 of
our dual solutions. The certificates turn this numerical agreement into exact
lower bounds; the 15 orbits settled by trees give dual value one, as they
must.

\emph{Lean formalization.} The Supplemental Material includes a
Lean~4/mathlib formalization of the exact algebraic and finite parts. Three
\texttt{native\_decide} checks reproduce the 35-task and 6435-task censuses of
decision trees and restricted trees, including the 395 cases settled by trees;
\texttt{native\_decide} also trusts Lean's compiler, through an extra axiom
that \texttt{\#print axioms} reports. The 43 dual certificates are checked by
the separate exact verifier, not in Lean. Density operators, general quantum
query algorithms, and the general parity lower bound are not formalized.

The census arranges balanced functions as answer partitions, identifies only
partitions related by symmetries of the oracle, and attaches a short exact
certificate to every class. Its proportions describe small cases only; they
are not a criterion for asymptotic speedup.

\section{Decohered access: Bernstein--Vazirani with an inaccessible witness}
\label{sec:witness}

For \(a,x\in\mathbb F_2^n\), let \(f_a(x)=a\cdot x \pmod2\) and
\(O_a\ket{x,y}=\ket{x,y\oplus a\cdot x}\). With ordinary phase kickback the
states for different \(a\) are orthogonal, so one query identifies
\(a\)~\cite{be-va}. We now let the oracle also leave an inaccessible record of
the address, without changing any value it returns on basis states.

Choose \(S\subseteq\{1,\ldots,n\}\) with \(|S|=k\), let \(x_S\) be the
corresponding substring of \(x\), and let \(\bar S\) be the complement of
\(S\). Each call acts as
\begin{equation}
 W_a^S\ket{x,y,e}
 =\ket{x,y\oplus a\cdot x,e\oplus x_S}
 \label{eq:witness-unitary}
\end{equation}
on the address \(x\), the target bit \(y\), and a fresh \(k\)-bit witness
register \(E\) with content \(e\), prepared in \(\ket{0^k}\) and then
discarded. The
oracle the algorithm sees is therefore the channel
\begin{equation}
 \Phi_a^S(\rho)=\operatorname{Tr}_E\!\left[
 W_a^S(\rho\otimes\ket{0^k}\!\bra{0^k})(W_a^S)^\dagger
 \right]
 =O_a\,\Delta_S(\rho)\,O_a^\dagger ,
 \label{eq:witness-channel}
\end{equation}
where \(\Delta_S(\rho)=\sum_s\Pi_s\rho\Pi_s\) dephases the address register in
the coordinates \(S\) (\(\Pi_s\) projects onto \(x_S=s\)); copying \(x_S\)
into a discarded register is exactly this dephasing, and \(O_a\) commutes with
every \(\Pi_s\). Basis queries still return \(a\cdot x\) exactly, but the
witness destroys coherence between addresses with different \(x_S\), as in
which-way experiments~\cite{Englert1996} and in the coherence analysis of the
Bernstein--Vazirani algorithm~\cite{NaseriEtAl2022}.

\emph{Surviving partition.} Define
\begin{equation}
 \Pi_S=\{C_u:u\in\mathbb F_2^{n-k}\},\qquad
 C_u=\{a:a_{\bar S}=u\}.
 \label{eq:surviving-partition}
\end{equation}
For the classical comparison, let \(\mathcal O^{\mathrm{val}}=\{O_a^{\mathrm{val}}:a\in
\mathbb F_2^n\}\) be the value oracle obtained by preparing and reading out
\(\Phi_a^S\) in the computational basis: a query chooses \(x\) and returns
\(a\cdot x\). For \(k<n\) and worst-case error \(\epsilon<1/2\),
\begin{equation}
 Q_\epsilon^{\Phi}(\Pi_S)=1,\qquad
 R_\epsilon(\Pi_S;\mathcal O^{\mathrm{val}})=n-k.
 \label{eq:partition-speedup}
\end{equation}
For the quantum part, the usual Bernstein--Vazirani input and phase kickback,
followed by Eq.~\eqref{eq:witness-channel}, leave on the \(n\)-qubit address
register (after factoring out the target qubit, which stays in \(\ket-\)) the
state
\begin{equation}
 \rho_a^S=2^{-n}\!\sum_{x,x':\,x_S=x'_S}
 (-1)^{a\cdot(x\oplus x')}\ket x\!\bra{x'}.
 \label{eq:reduced-bv}
\end{equation}
After the final Hadamard gates, the orthogonality of characters gives, with
the coordinates reordered and \(z\) the measured address string,
\begin{equation}
\begin{aligned}
 H^{\otimes n}\rho_a^S H^{\otimes n}
 &=\ket{a_{\bar S}}\!\bra{a_{\bar S}}\otimes\frac{I_S}{2^k},\\
 \Pr(z|a)&=2^{-k}\,\mathbf 1[z_{\bar S}=a_{\bar S}].
\end{aligned}
 \label{eq:leaky-distribution}
\end{equation}
So one query resolves \(\Pi_S\) exactly; zero queries cannot, since the output
would not depend on the cell. Classically, the \(n-k\) basis queries outside
\(S\) suffice. For the classical lower bound, grant even the promise
\(a_S=0\) and let \(a_{\bar S}\) be uniformly distributed. Fewer than \(n-k\)
linear equations leave at least two equally likely candidates, so a decoder
that sees the whole transcript succeeds with probability at most \(1/2\).
Averaging over the protocol's random choices does not change this, so some
input has success at most \(1/2\).

\emph{Identifying \(a\).} The same channel also changes the cost of the
original task, identifying \(a\), that is, resolving the partition
\(\Pi_{\rm single}\) into single instances.

\begin{proposition}[Witness cost of identification]\label{prop:witness}
For \(0\le k\le n\), with \(S\) known to the algorithm, the exact quantum cost
of identifying \(a\) through \(\{\Phi_a^S\}\) is
\(Q_E^{\Phi}(\Pi_{\rm single})=\min\{k+1,n\}\); classically,
\(R_\epsilon(\Pi_{\rm single};\mathcal O^{\mathrm{val}})=n\) for
\(\epsilon<1/2\).
\end{proposition}

\begin{proof}
Upper bound: one Bernstein--Vazirani query with \(x_S=0\) returns
\(a_{\bar S}\), and \(k\) basis queries \(x=e_i\), \(i\in S\), return
\(a_S\); alternatively, \(n\) basis queries return \(a\). At \(k=0\) this is
the ordinary one-query algorithm. Lower bound for \(k\ge1\): record each
dephasing in Eq.~\eqref{eq:witness-channel} in an inaccessible environment
register. For a record \(\vec s=(s_1,\ldots,s_T)\), \(s_t\in\mathbb F_2^k\), let
\(\ket{\phi_a(\vec s)}\) be the conditional unnormalized output, so that
\(\rho_a=\sum_{\vec s}\ket{\phi_a(\vec s)}\!\bra{\phi_a(\vec s)}\). Because the
records are inaccessible, exact identification requires
\(\braket{\phi_a(\vec s)|\phi_{a'}(\vec s')}=0\) for all \(a\neq a'\) and all
records \(\vec s,\vec s'\). In the target qubit's \(X\) basis, the \(t\)th query
multiplies the component at address \(x=(s_t,z_t)\), where \(z_t=x_{\bar S}\)
holds the unrecorded coordinates, either by 1 or by
\((-1)^{a\cdot(s_t,z_t)}\). Hence \(\phi_a(\vec s)\) is a linear combination of
characters \((-1)^{a\cdot f}\) whose frequencies are \(f=0\) or
\(f=\sum_{t\in F}(s_t,z_t)\) for nonempty subsets \(F\) of the queries.
(i) If \(s_1,\ldots,s_T\) do not span \(\mathbb F_2^k\), choose \(h\neq0\)
orthogonal to all \(s_t\). Every such character is then invariant under
\(a\mapsto a+(h,0)\), so \(\phi_{a+(h,0)}(\vec s)=\phi_a(\vec s)\), and
orthogonality forces \(\phi_a(\vec s)=0\). Hence \(T\ge k\), and only records
whose \(s_t\) span \(\mathbb F_2^k\) contribute. (ii) Let \(T=k<n\). In every
contributing record the \(s_t\) are linearly independent, so every frequency
with \(F\neq\emptyset\) has a nonzero \(S\) part, and for \(u\neq0\) the sum
\(\sum_a(-1)^{u\cdot a_{\bar S}}\phi_a(\vec s)\) vanishes. Since the
\(\phi_a(\vec s)\) are pairwise orthogonal, the squared norm of this sum is
\(\sum_a\lVert\phi_a(\vec s)\rVert^2\), so \(\phi_a(\vec s)=0\) for every
contributing record, contradicting normalization. Hence \(T\ge k+1\) for
\(k<n\). The classical bound is the singleton case of the argument above.
\end{proof}

At \(k=0\) nothing is recorded, and Eq.~\eqref{eq:leaky-distribution} is the
ordinary Bernstein--Vazirani output, with one quantum query against \(n\)
classical ones. Each recorded address bit adds exactly one query until the
cost reaches the classical value \(n\); from \(k=n-1\) on, the quantum
advantage for identification is gone, while the classical cost stays \(n\).
If the witness were available to the algorithm, the fixed coupling
\(\ket{x,e}\mapsto\ket{x,e\oplus x_S}\) could be undone without another oracle
call, restoring the ordinary output. What the algorithm accesses here is a
channel, not the unitary \(W_a^S\): as with in-place versus standard
oracles~\cite{KashefiKentVedralBanaszek2002} and faulty
oracles~\cite{RegevSchiff2008}, whether the oracle leaves an inaccessible
record is part of the problem.

\emph{Partially distinguishable witnesses.} Suppose the two witness states
\(\ket{\xi_0^{(i)}},\ket{\xi_1^{(i)}}\) of a recorded coordinate \(i\in S\) are not orthogonal but have real overlap
\(\kappa_i=\braket{\xi_0^{(i)}|\xi_1^{(i)}}\in[0,1)\). Each call then
dephases that coordinate only partially,
\(\rho\mapsto\frac{1+\kappa_i}{2}\rho+\frac{1-\kappa_i}{2}Z_i\rho Z_i\), with
\(Z_i\) the Pauli \(Z\) on address bit \(i\); orthogonal witnesses,
\(\kappa_i=0\), give the dephasing of Eq.~\eqref{eq:witness-channel}. One
Bernstein--Vazirani query returns \(z=a\) with probability
\(\prod_{i\in S}(1+\kappa_i)/2\), the which-way visibility of
Ref.~\cite{Englert1996}. A single call therefore still identifies \(a\) with
bounded error whenever \(\prod_{i\in S}(1+\kappa_i)>2^{k-1}\), for instance
for every \(\kappa_i>0\) when \(k=1\). Exactness, however, is lost for every
\(\kappa_i<1\). Both terms of the partial dephasing carry positive weight, so
after any number of calls the output states have the same supports as for
\(\kappa_i=0\), and exact identification faces the same orthogonality
conditions; Proposition~\ref{prop:witness} holds unchanged. Only
\(\kappa_i=1\), no leakage, restores identification in one query.

\section{Discussion}
\label{sec:discussion}

Can the comparison go the other way, with a classical algorithm needing fewer
oracle calls than every quantum one? Not in the standard models of
computation. A quantum computer can run any classical algorithm: reversible
classical gates such as the Toffoli gate are unitary, and measurements supply
random bits. Bounded-error quantum computation therefore contains
bounded-error classical computation, also relative to oracles~\cite{be-va},
and the exact quantum query complexity of a function never exceeds its
deterministic query complexity. The same holds for the access models of this
paper. Each either returns the classical value on a basis-state query, so
that a quantum algorithm can run the classical protocol call for call, or,
like the sign oracle of Proposition~\ref{prop:sign}, returns nothing a
classical protocol could use. What can happen is that the advantage
disappears. The qutrit cyclic shift needs two calls for Deutsch's problem,
as many as a classical algorithm (Corollary~\ref{cor:completions}); the
parity of \(N\) bits needs \(N/2\) quantum queries~\cite{Farhi-98,bbcmw-01};
Grover's quadratic speedup for unstructured search is
optimal~\cite{benn:97}; and for total Boolean functions, quantum and
deterministic query complexities are polynomially related~\cite{bbcmw-01}.

All results of this paper follow one strategy, which can be stated in four
steps; the physics enters in the middle two.

\begin{itemize}
\item
\emph{Answer partition \(\Pi_g\)} [Eq.~\eqref{eq:task-partition}]:
Put two hidden instances in the same cell whenever the task has the same
correct answer for both. The cells \(C_j\) are sets of classical instances,
not subspaces. The algorithm must learn the cell, and only the cell: in
Deutsch's problem the four possible functions fall into two cells,
constant and balanced, and the algorithm never has to find out which
function within a cell it faces.

\item
\emph{Oracle access} (Secs.~\ref{sec:spectrum} and~\ref{sec:witness}):
Fix what one oracle call does. The partition fixes the correct answers to
the question; the access model fixes how that answer can be obtained by
quantum state discrimination.

\item
\emph{Probe and one quantum oracle call} (Sec.~\ref{sec:criterion}):
Choose an instance-independent probe state, for instance a uniform
superposition obtained by ``fanning out'' a single basis state with
Hadamard gates. One oracle call turns the probe into the output state
\(\ket{\Psi_a}\) of the hidden instance \(a\).

\item
\emph{Orthogonality across cells} (Proposition~\ref{prop:block}):
Ask whether some probe makes \(\braket{\Psi_a|\Psi_{a'}}=0\) whenever
\(g(a)\ne g(a')\) [Eq.~\eqref{eq:criterion}]. A negative answer needs no
search over the infinitely many probes: a single matrix, a dual certificate
(Proposition~\ref{prop:dual}), proves that none works. Only orthogonal
states can be told apart with certainty, so exactness requires
orthogonality, but only between cells. Within a cell, states may overlap or
even coincide: the measurement reads the cell label, not the instance, as a
post office sorts letters by postal code without reading the street.
\end{itemize}

Two equally weighted branches of a quantum state cancel, that is, interfere
destructively, only if one carries an extra factor \(-1\) relative to the
other. In the examples of the Introduction, the standard flip on a qubit
and the transposition on a qutrit both have the eigenvalue \(-1\) and
separate constant from balanced functions in one call; the qutrit cyclic
shift, whose eigenvalues are the cube roots of unity, cannot
(Theorem~\ref{thm:spectrum} and Corollary~\ref{cor:completions}). With
more than two branches, other phases can cancel, for instance
\(1+\omega+\omega^2=0\) for \(\omega=e^{2\pi i/3}\); this is the interference
used by the quantum Fourier transform and Shor's algorithm. Deutsch's
problem, however, compares only two address branches, and therefore needs
\(-1\). With more calls, the same orthogonality is required of the final states, with
instance-independent operations between the calls (Sec.~\ref{sec:census});
with bounded error, it is relaxed to a required success probability
(Proposition~\ref{prop:bounded}).

Some of these separations need no non-Clifford operations. The standard
Deutsch and Bernstein--Vazirani protocols, including their XOR oracles, use
only Clifford operations, and so does the witness coupling of
Eq.~\eqref{eq:witness-unitary} before the witness is discarded. Given a full
classical description of the circuit, they can therefore be simulated
efficiently~\cite{GottesmanKnill,AaronsonGottesman2004}. In the query model,
however, such a description would reveal the hidden oracle data, which is more
than calling the oracle provides. For these examples, nonstabilizer resources
(``magic''~\cite{Veitch2014}) are not needed for a query advantage. General
oracles of this kind with a binary-encoded address need multiply controlled gates,
which are not Clifford, so this remark is specific to Deutsch and
Bernstein--Vazirani.

The results vary these steps one at a time. Corollary~\ref{cor:completions}
and Proposition~\ref{prop:witness} change the access (the second step) while
holding the task and its classical values fixed, and the exact cost
changes; Proposition~\ref{prop:bounded} shows what survives with bounded
error. Theorem~\ref{thm:spectrum} and Proposition~\ref{prop:oneclass} give the
rule for every response unitary. The sign oracle of
Proposition~\ref{prop:sign} is a deliberately weaker oracle, whose
values are invisible to basis queries. The census keeps the standard oracle
fixed and varies the balanced task (the first step), and the cost changes
again; for \(N\le4\) every class is certified exactly. Shor's and Grover's
algorithms~\cite{shor:94,grover} fit the same picture: each exploits structure
of the oracle family that matches a specific coherent access.

Partitions are therefore needed to state a task but do not decide its
complexity, and neither do classical values. A quantum advantage occurs
when interference matched to the oracle's structure separates the answer cells
in fewer calls than any classical protocol needs with the same access.
Superlinear exact separations for total Boolean functions are
known~\cite{Ambainis2016SuperlinearExact}, but an asymptotic speedup requires
a structured family of task--oracle pairs; percentages from a finite census
cannot replace that analysis. Within the present framework, the parity
question for the qutrit shift (Sec.~\ref{sec:spectrum}) is a natural test of
whether a response-spectrum effect persists for all \(N\). A second open
question is which response \(V\) is optimal for several queries. A response
with eigenvalues \(1\) and \(-1\) reproduces the XOR oracle exactly in each
call, by encoding the control qubit in the corresponding eigenvectors, and the
degree bound holds for every \(V\); whether some other response beats the XOR
oracle for some task with two or more queries remains open.

In summary, the paper adds three things. First, a complete one-query
classification for every response unitary (Theorem~\ref{thm:spectrum} and
Proposition~\ref{prop:oneclass}), which extends the XOR-oracle
characterization of Chen, Ye, and Li~\cite{ChenYeLi2020}. Second, a caution
for implementations: ``a quantum oracle for \(f\)'' is underspecified, since
reversible oracles with the same classical readout can differ in quantum
cost. This matters in particular for qudit realizations, where the natural
reversible extension may be a cyclic shift without the eigenvalue \(-1\); the
one-query solution of Deutsch's problem is then lost. Third, exact results
where only numerical ones existed: rational certificates for the 6040
four-bit cases, the optimal bounded-error success of odd cyclic responses,
and the exact cost \(\min\{k+1,n\}\) of identification with an inaccessible
address copy. All effects found here are constant factors; none is an
asymptotic speedup.

\begin{acknowledgments}
The author thanks Noson S. Yanofsky for discussions that motivated this
work. Two AI assistants were used. OpenAI Codex ChatGPT
Astra ultra
supported structural editing, consistency checks, the implementation of the
census code, and the Lean~4 formalization. Claude Opus~5.5 (Anthropic) was
used in the revision to review the manuscript and the literature, to draft
revised text---including the statements and proofs of
Propositions~\ref{prop:dual}--\ref{prop:witness} and the new proof of
Theorem~\ref{thm:spectrum}---and to carry out the bounded-error and
dual-certificate computations of Secs.~\ref{sec:spectrum} and~\ref{sec:census}.
The author directed both tools, checked the mathematical arguments, and
verified the computations as described in Sec.~\ref{sec:census}; the author
accepts full responsibility for the content and code. This research
was funded in whole or in part by the Austrian Science Fund (FWF), grant
\href{https://doi.org/10.55776/PIN5424624}{PIN5424624}. The author
acknowledges TU Wien Bibliothek for support through its Open Access Funding
Programme.
\end{acknowledgments}

\noindent\textit{Data and code availability.---}See Supplemental Material at
[URL will be inserted by publisher] for the Python and Lean~4 sources, the
proof report, the 43 rational dual certificates with an exact verifier, the
census tables and witness files, and the cross-validation against the
published NPN data of Ref.~\cite{MontanaroJozsaMitchison2015}.

\bibliography{svozil}

\end{document}
